% !TeX root = ../sections/02_exercises.tex

\begin{exercise}
	環
	\[
	R=\mathbb{Z}[\sqrt{2},\sqrt{3}]
	\]
	とそのイデアル
	\[
	I=(\sqrt{6})
	\]
	について考える。
	\(\mathbb{Z}\)-加群としての同型写像
	\[
	f:\mathbb{Z}[\sqrt{2},\sqrt{3}]
	\longrightarrow
	\mathbb{Z}^4
	\]
	を
	\[
	f(a+b\sqrt{2}+c\sqrt{3}+d\sqrt{6})
	=
	\begin{pmatrix}
		a\\
		b\\
		c\\
		d
	\end{pmatrix}
	\]
	で定める。
	\(V=f(I)\) を，自由 \(\mathbb{Z}\)-加群 \(\mathbb{Z}^4\) の
	\(\mathbb{Z}\)-部分加群とする。
	
	\begin{enumerate}
		\item
		\(V\) の階数 \(m\) および \(V\) の基底
		\[
		[a_1,\ldots,a_m]
		\]
		を求めよ。
		
		\item
		\(A=(a_1,\ldots,a_m)\) とおく。
		\(A\) の単因子を求めよ。
		
		\item
		\(R/I\) を \(\mathbb{Z}\)-加群と見たときの不変系を求めよ。
	\end{enumerate}
\end{exercise}

\begin{background}
	この問題では，環
	\[
	R=\mathbb{Z}[\sqrt{2},\sqrt{3}]
	\]
	を \(\mathbb{Z}\)-加群として見る。
	任意の元は
	\[
	a+b\sqrt{2}+c\sqrt{3}+d\sqrt{6}
	\qquad
	(a,b,c,d\in\mathbb{Z})
	\]
	と一意的に表されるので，
	\(R\) は \(\mathbb{Z}\)-加群として
	\[
	R\simeq \mathbb{Z}^4
	\]
	である。
	
	この同型が問題文の写像
	\[
	f:R\to\mathbb{Z}^4
	\]
	である。
	すなわち，
	\[
	f(a+b\sqrt{2}+c\sqrt{3}+d\sqrt{6})
	=
	\begin{pmatrix}
		a\\
		b\\
		c\\
		d
	\end{pmatrix}
	\]
	によって，\(R\) の元を係数ベクトルに対応させる。
	
	イデアル
	\[
	I=(\sqrt{6})
	\]
	は，\(R\) の元に \(\sqrt{6}\) を掛けたもの全体である。
	したがって，
	\[
	I=
	\left\{
	(a+b\sqrt{2}+c\sqrt{3}+d\sqrt{6})\sqrt{6}
	\ \middle|\ 
	a,b,c,d\in\mathbb{Z}
	\right\}
	\]
	である。
	
	積を展開すると，
	\[
	\sqrt{6}\cdot 1=\sqrt{6},
	\qquad
	\sqrt{6}\cdot \sqrt{2}=2\sqrt{3},
	\]
	\[
	\sqrt{6}\cdot \sqrt{3}=3\sqrt{2},
	\qquad
	\sqrt{6}\cdot \sqrt{6}=6
	\]
	となる。
	これにより，\(I\) の生成元を \(\mathbb{Z}^4\) のベクトルとして表せる。
	
	剰余 \(R/I\) を \(\mathbb{Z}\)-加群として調べることは，
	同型 \(f\) によって
	\[
	\mathbb{Z}^4/V
	\]
	を調べることに対応する。
	したがって，\(V=f(I)\) の基底を取り，
	その基底を列に並べた行列 \(A\) の単因子を求めればよい。
\end{background}

\begin{strategy}
	まず，イデアル \(I=(\sqrt{6})\) の元を具体的に展開する。
	任意の \(R\) の元は
	\[
	a+b\sqrt{2}+c\sqrt{3}+d\sqrt{6}
	\]
	と書けるので，
	\[
	(a+b\sqrt{2}+c\sqrt{3}+d\sqrt{6})\sqrt{6}
	\]
	を計算する。
	
	展開すると，
	\[
	a\sqrt{6}
	+
	b\sqrt{12}
	+
	c\sqrt{18}
	+
	d\sqrt{36}
	\]
	である。
	ここで
	\[
	\sqrt{12}=2\sqrt{3},
	\qquad
	\sqrt{18}=3\sqrt{2},
	\qquad
	\sqrt{36}=6
	\]
	だから，
	\[
	(a+b\sqrt{2}+c\sqrt{3}+d\sqrt{6})\sqrt{6}
	=
	6d+3c\sqrt{2}+2b\sqrt{3}+a\sqrt{6}
	\]
	となる。
	
	\begin{enumerate}
		\item[(1)]
		上の式を \(f\) によって \(\mathbb{Z}^4\) のベクトルに移す。
		すると，
		\[
		f(6)=
		\begin{pmatrix}
			6\\
			0\\
			0\\
			0
		\end{pmatrix},
		\qquad
		f(3\sqrt{2})=
		\begin{pmatrix}
			0\\
			3\\
			0\\
			0
		\end{pmatrix},
		\]
		\[
		f(2\sqrt{3})=
		\begin{pmatrix}
			0\\
			0\\
			2\\
			0
		\end{pmatrix},
		\qquad
		f(\sqrt{6})=
		\begin{pmatrix}
			0\\
			0\\
			0\\
			1
		\end{pmatrix}
		\]
		が得られる。
		
		したがって，
		\[
		V=f(I)
		\]
		はこれらのベクトルによって生成される。
		これらが \(\mathbb{Z}\)-一次独立であることを確認すれば，
		\(V\) の階数と基底が得られる。
		
		\item[(2)]
		(1) で得た基底を列ベクトルとして並べ，
		\[
		A=(a_1,\ldots,a_m)
		\]
		を作る。
		その行列に対して整数係数の行基本変形・列基本変形を行い，
		単因子標準形に直す。
		
		単因子を求めるときは，
		小行列式の最大公約数を使って確認するとよい。
		すなわち，
		\[
		\Delta_k=
		\gcd\{\text{\(A\) の \(k\) 次小行列式}\}
		\]
		を求める。
		そのとき，
		\[
		d_1=\Delta_1,
		\qquad
		d_2=\frac{\Delta_2}{\Delta_1},
		\qquad
		d_3=\frac{\Delta_3}{\Delta_2},
		\qquad
		d_4=\frac{\Delta_4}{\Delta_3}
		\]
		として単因子が得られる。
		
		\item[(3)]
		\(R/I\) を \(\mathbb{Z}\)-加群として見ると，
		同型 \(f\) によって
		\[
		R/I\simeq \mathbb{Z}^4/V
		\]
		と考えられる。
		
		したがって，(2) で求めた単因子
		\[
		d_1,d_2,d_3,d_4
		\]
		を用いて，
		\[
		R/I
		\simeq
		\mathbb{Z}/d_1\mathbb{Z}
		\oplus
		\mathbb{Z}/d_2\mathbb{Z}
		\oplus
		\mathbb{Z}/d_3\mathbb{Z}
		\oplus
		\mathbb{Z}/d_4\mathbb{Z}
		\]
		と表す。
		
		ただし，\(d_i=1\) の成分は自明な群であるため，
		不変系を書くときには通常省略する。
	\end{enumerate}
\end{strategy}